\documentclass{article}

\usepackage{amsmath}
\usepackage[margin=1in]{geometry}

\begin{document}

We want to find an expression for the light-cone momentum and virtuality distribution, $\tilde{\rho}(\alpha,v)$ for deuterium.
  Segarra et al., Eq.~5, gives us:
  \begin{equation}
    \tilde{\rho}(\alpha,v) = \int dE d^3p S(E,\vec{p}) \frac{E+p_z}{E} \delta\left( \alpha - \frac{A(E+p_z)}{m_A}\right)
    \delta\left( v - \frac{E^2 - |p|^2 - m_N^2}{m_N^2}\right).
    \label{eq:setup}
  \end{equation}

  For deuterium, where the $(A-1)$ system is merely a nucleon, the spectral function is simply a momentum density, $n(p)$, with an
  energy-conserving $\delta$-function:
  \begin{equation}
    S^d(E,p) = n(p) \delta\left(E - m_d + \sqrt{m_N^2 + p^2}\right).
  \end{equation}

  For the derivation that follows, I will preserve the distinction between the proton and neutron mass. The nucleon is struck
  will be termed the ``leading'' nucleon, with subscript `L,' while the spectator nucleon will receive subscript `s.' In this
  way, the spectral function for striking a nucleon of species `L' can be written:
  \begin{equation}
    S_L^d(E,p) = n(p) \delta\left(E - m_d + \sqrt{m_s^2 + p^2}\right),
  \end{equation}
  where the mass of spectator species enters into the $\delta$-function.

  Plugging this spectral function into equation~\ref{eq:setup} gives us:
  \begin{equation}
    \tilde{\rho}_L^d(\alpha,v) = \int dE d^3p n(p) \delta\left(E - m_d + \sqrt{m_s^2 + p^2}\right) \frac{E+p_z}{E}
    \delta\left( \alpha - \frac{2(E+p_z)}{m_d}\right)
    \delta\left( v - \frac{E^2 - |p|^2 - m_L^2}{m_L^2}\right).
  \end{equation}
  It is simplest to attack this in cylindrical coordinates: $d^3p \rightarrow dp_z p_\perp dp_\perp d\phi$. We can eliminate
  the $p_\perp$ integral using the virtuality $\delta$-function, from which we can see that $p_\perp = \sqrt{E^2 -p_z^2 + m_L^2(1+v)}$.
  This yields:
  \begin{equation}
    \begin{split}
    \tilde{\rho}_L^d(\alpha,v) = \int dE dp_z d\phi p_\perp dp_\perp n(p) \delta\left(E - m_d + \sqrt{m_s^2 + p^2}\right) \frac{E+p_z}{E}
    \delta\left( \alpha - \frac{2(E+p_z)}{m_d}\right)\\
    \cdot \frac{m_L^2}{2p_\perp} \delta\left( p_\perp - \sqrt{E^2 -p_z^2 + m_L^2(1+v)}\right),
    \end{split}
  \end{equation}
  which simplifies to
  \begin{equation}
    \tilde{\rho}_L^d(\alpha,v) = \frac{m_L^2}{2} \int dE dp_z d\phi n(p) \delta\left(E - m_d + \sqrt{m_s^2 + p^2}\right) \frac{E+p_z}{E}
    \delta\left( \alpha - \frac{2(E+p_z)}{m_d}\right).
  \end{equation}
  The momentum magnitude, $p$, can be expressed as $p^2 = p_\perp^2 + p_z^2$, which transforms the penultimate delta function:
  \begin{equation}
    \tilde{\rho}_L^d(\alpha,v) = \frac{m_L^2}{2} \int dE dp_z d\phi n(p) \delta\left(E - m_d + \sqrt{m_s^2 + E^2 - m_L^2(1+v)}\right) \frac{E+p_z}{E}
    \delta\left( \alpha - \frac{2(E+p_z)}{m_d}\right).
  \end{equation}
  There is no sign of any $\phi$ dependence, so we may remove that integral,
  \begin{equation}
    \tilde{\rho}_L^d(\alpha,v) = \pi m_L^2 \int dE dp_z n(p) \delta\left(E - m_d + \sqrt{m_s^2 + E^2 - m_L^2(1+v)}\right) \frac{E+p_z}{E}
    \delta\left( \alpha - \frac{2(E+p_z)}{m_d}\right),
  \end{equation}
  and turn our attention to the $\delta$-function specifying $\alpha$. We can see that $p_z = \frac{m_d \alpha}{2} - E$, which allows us to write:
  \begin{equation}
    \tilde{\rho}_L^d(\alpha,v) = \pi m_L^2 \int dE dp_z n(p) \delta\left(E - m_d + \sqrt{m_s^2 + E^2 - m_L^2(1+v)}\right) \frac{E+p_z}{E}
    \frac{m_d}{2} \delta\left( p_z - \frac{m_d \alpha}{2} + E \right),
  \end{equation}
  perform the $dp_z$ integral, and end up with:
  \begin{equation}
    \tilde{\rho}_L^d(\alpha,v) = \frac{\pi m_d^2 m_L^2}{4} \int \frac{dE}{E} n(p) \delta\left(E - m_d + \sqrt{m_s^2 + E^2 - m_L^2(1+v)}\right).
  \end{equation}


  To tackle the $dE$ integral, we must transform the energy $\delta$-function. Treating it as $\delta(f(E)$, where
  $f(E) = E - m_d + \sqrt{m_s^2 + E^2 - m_L^2(1+v)}$, we can take the derivative of $f$ with respect to $E$:
  \begin{equation}
    f'(E) = 1 + \frac{E}{\sqrt{E^2 + m_s^2 - m_L^2(1+v)}}.
  \end{equation}
  We can also invert $f$ to get
  \begin{equation}
    E= \frac{m_d^2 + m_L^2(1+v) - m_s^2}{2m_d}. 
  \end{equation}
  This allows us to transform the energy $\delta$-function to get:
  \begin{equation}
    \tilde{\rho}_L^d(\alpha,v) = \frac{\pi m_d^2 m_L^2}{4} \int \frac{dE}{E} n(p) \left|
    \frac{\sqrt{E^2 + m_s^2 - m_L^2(1+v)}}{E+\sqrt{E^2 + m_s^2 - m_L^2(1+v)}} \right|
    \delta\left(E -\frac{m_d^2 + m_L^2(1+v) - m_s^2}{2m_d}\right).
  \end{equation}
  Performing the integral gives us:
  \begin{equation}
    \tilde{\rho}_L^d(\alpha,v) = \frac{\pi m_d^2 m_L^2}{4}\cdot \frac{n(p)}{E} \cdot \left|
    \frac{\sqrt{E^2 + m_s^2 - m_L^2(1+v)}}{E+\sqrt{E^2 + m_s^2 - m_L^2(1+v)}} \right|,
  \end{equation}
  and we must now substitute in $\frac{m_d^2 + m_L^2(1+v) - m_s^2}{2m_d}$ for $E$.

  Let's first attack the term in the square root:
  \begin{align}
    E^2 &= \frac{1}{4m_d^2} \left[
      m_d^4 + 2m_d^2m_L^2(1+v) - 2m_d^2m_s^2 + m_L^4(1+v)^2 - 2m_L^2m_s^2(1+v) +m_s^4\right] \\
    m_s^2 - m_L^2(1+v) + E^2 &= \frac{1}{4m_d^2} \left[
      m_d^4 - 2m_d^2m_L^2(1+v) + 2m_d^2m_s^2 + m_L^4(1+v)^2 - 2m_L^2m_s^2(1+v) +m_s^4\right] \\
    m_s^2 - m_L^2(1+v) + E^2 &= \left( \frac{1}{2m_d} \left[
      m_d^2 - m_L^2(1+v) + m_s^2\right]\right)^2\\
    \sqrt{m_s^2 - m_L^2(1+v) + E^2} &= \frac{\left| m_d^2 - m_L^2(1+v) + m_s^2\right|}{2m_d}
  \end{align}

  Substituting this result into our expression gives us:
  \begin{align}
    \tilde{\rho}_L^d(\alpha,v) &= \frac{\pi m_d^2 m_L^2}{4} \cdot n(p) \cdot \frac{2m_d}{m_d^2 + m_L^2(1+v) - m_s^2} \cdot
    \left|
    \frac{m_d^2 - m_L^2(1+v) + m_s^2}{m_d^2 - m_L^2(1+v) + m_s^2 + m_d^2 + m_L^2(1+v) - m_s^2} \right|\\
    &= \frac{\pi m_d^2 m_L^2}{4} \cdot n(p) \cdot \frac{2m_d}{m_d^2 + m_L^2(1+v) - m_s^2} \cdot
    \left|
    \frac{m_d^2 - m_L^2(1+v) + m_s^2}{2 m_d^2 } \right|\\
    &= \frac{\pi m_d m_L^2}{4} \cdot n(p) \cdot \left| \frac{m_d^2 - m_L^2(1+v) + m_s^2}{m_d^2 + m_L^2(1+v) - m_s^2} \right|   
  \end{align}

  Up until now, we have preserved the distinction between the proton mass, $m_p$, and the neutron mass, $m_n$. However, if
  ignore that distinction and assign both with the same nucleon mass, $m_N = m_s = m_L$, then our formula reduces to:
  \begin{equation}
    \tilde{\rho}^d(\alpha,v) = \frac{\pi m_d m_N^2}{4} \cdot n(p) \cdot \left| \frac{m_d^2 - vm_N^2}{m_d^2 + vm_N^2} \right|   
  \end{equation}
\end{document}
